\documentclass[10pt,a4paper]{amsart}
\usepackage{amsmath,amssymb,amsthm,graphicx}
% For commutative diagrams we use Xy-pic
\usepackage[all]{xy}

% We use 2cell for 2-commutative diagrams.
\xyoption{2cell}
% Theorem environments.
%
\theoremstyle{plain}
\newtheorem{theorem}[subsection]{Theorem}
\newtheorem{proposition}[subsection]{Proposition}
\newtheorem{lemma}[subsection]{Lemma}

\theoremstyle{definition}
\newtheorem{definition}[subsection]{Definition}
\newtheorem{example}[subsection]{Example}
\newtheorem{exercise}[subsection]{Exercise}
\newtheorem{situation}[subsection]{Situation}

\theoremstyle{remark}
\newtheorem{remark}[subsection]{Remark}
\newtheorem{remarks}[subsection]{Remarks}

\numberwithin{equation}{subsection}

% Macros
%
\def\lim{\mathop{\mathrm{lim}}\nolimits}
\def\colim{\mathop{\mathrm{colim}}\nolimits}
\def\Spec{\mathop{\mathrm{Spec}}}
\def\Hom{\mathop{\mathrm{Hom}}\nolimits}
\def\Ext{\mathop{\mathrm{Ext}}\nolimits}
\def\SheafHom{\mathop{\mathcal{H}\!\mathit{om}}\nolimits}
\def\SheafExt{\mathop{\mathcal{E}\!\mathit{xt}}\nolimits}
\def\Sch{\mathit{Sch}}
\def\Mor{\mathop{\mathrm{Mor}}\nolimits}
\def\Ob{\mathop{\mathrm{Ob}}\nolimits}
\def\Sh{\mathop{\mathit{Sh}}\nolimits}
\def\NL{\mathop{N\!L}\nolimits}
\def\CH{\mathop{\mathrm{CH}}\nolimits}
\def\proetale{{pro\text{-}\acute{e}tale}}
\def\etale{{\acute{e}tale}}
\def\QCoh{\mathit{QCoh}}
\def\Ker{\mathop{\mathrm{Ker}}}
\def\Im{\mathop{\mathrm{Im}}}
\def\Coker{\mathop{\mathrm{Coker}}}
\def\Coim{\mathop{\mathrm{Coim}}}

% Boxtimes
%
\DeclareMathSymbol{\boxtimes}{\mathbin}{AMSa}{"02}

%
% Macros for moduli stacks/spaces
%
\def\QCohstack{\mathcal{QC}\!\mathit{oh}}
\def\Cohstack{\mathcal{C}\!\mathit{oh}}
\def\Spacesstack{\mathcal{S}\!\mathit{paces}}
\def\Quotfunctor{\mathrm{Quot}}
\def\Hilbfunctor{\mathrm{Hilb}}
\def\Curvesstack{\mathcal{C}\!\mathit{urves}}
\def\Polarizedstack{\mathcal{P}\!\mathit{olarized}}
\def\Complexesstack{\mathcal{C}\!\mathit{omplexes}}
% \Pic is the operator that assigns to X its picard group, usage \Pic(X)
% \Picardstack_{X/B} denotes the Picard stack of X over B
% \Picardfunctor_{X/B} denotes the Picard functor of X over B
\def\Pic{\mathop{\mathrm{Pic}}\nolimits}
\def\Picardstack{\mathcal{P}\!\mathit{ic}}
\def\Picardfunctor{\mathrm{Pic}}
\def\Deformationcategory{\mathcal{D}\!\mathit{ef}}
\usepackage{cleveref}
\crefname{originalsection}{original section}{original sections}
\Crefname{originalsection}{Original section}{Original sections}
\makeatletter
\newlabel{derived-section-canonical-delta-functor}{{014Z}{source}{}{sourcepointer-1}{}}
\newlabel{derived-section-canonical-delta-functor@cref}{{[originalsection][0][]014Z}{[1][0][]source}{}{}{}}
\makeatother
\crefname{originallemma}{original lemma}{original lemmas}
\Crefname{originallemma}{Original lemma}{Original lemmas}
\makeatletter
\newlabel{lemma-proper-base-change}{{09V5}{source}{}{sourcepointer-2}{}}
\newlabel{lemma-proper-base-change@cref}{{[originallemma][0][]09V5}{[1][0][]source}{}{}{}}
\makeatother
\crefname{originallemma}{original lemma}{original lemmas}
\Crefname{originallemma}{Original lemma}{Original lemmas}
\makeatletter
\newlabel{lemma-proper-base-change-f-star}{{0A3U}{source}{}{sourcepointer-3}{}}
\newlabel{lemma-proper-base-change-f-star@cref}{{[originallemma][0][]0A3U}{[1][0][]source}{}{}{}}
\makeatother
\crefname{originalproposition}{original proposition}{original propositions}
\Crefname{originalproposition}{Original proposition}{Original propositions}
\makeatletter
\newlabel{proposition-finite-higher-direct-image-zero}{{03QP}{source}{}{sourcepointer-4}{}}
\newlabel{proposition-finite-higher-direct-image-zero@cref}{{[originalproposition][0][]03QP}{[1][0][]source}{}{}{}}
\makeatother
\title{The blow-up square of a finite type quasi-coherent ideal gives a cohomological descent distinguished triangle for torsion coefficient complexes}
\author{The Stacks Project Authors}
\date{stacks-project:etale-cohomology.tex:local-snapshot-20261011; GNU FDL 1.2 or later}
\begin{document}
\maketitle
\noindent\textit{Editorial note.} All delivered mathematical text is the author's own, copied from the pinned Stacks Project chapter file etale-cohomology.tex; no statement or proof step was rewritten, reordered or removed.\par\medskip
\noindent\textit{Editorial note.} The counted proof is the author's complete argument: the proof is the author's complete argument: the choice of the complexes representing the three derived pushforwards, the identification by proper base change, the reduction to a stalk statement via the triangle attached to a short exact sequence of complexes, and the two-case stalk verification are all written out.\par\medskip
\noindent\textit{Editorial note.} The label lemma-blow-up-square-cohomological-descent names the counted statement, and the proof environment immediately following it in the source is delivered verbatim as its proof.\par\medskip
\noindent\textit{Editorial note.} The statement and proof are copied from the pinned chapter file; the four cited inputs are kept as external pointers with their canonical Stacks tags.\par\medskip
\noindent\textit{Editorial note.} The preamble is the chapter's own native preamble excerpt already used for this source: the Xy-pic setup (native preamble lines 14-18) and the native theorem declarations with the mathematical macros (native preamble lines 154-217).\par\medskip
\noindent\textit{Allowed standard prerequisites (editorial).}\begin{itemize}
\item The blow-up square of a closed subscheme cut out by a quasi-coherent ideal of finite type. Same chapter, the blow-up square with E = Z x\_X X' and the equalities c = i o pi = b o j. Supplies the geometric set-up of the statement and the notation for the four morphisms.
\item Proper base change and its pushforward form for the blow-up. Same chapter, the proper base change theorem and its consequence identifying pi\_*(I\textasciicircum{}bullet|\_E) with (b\_*I\textasciicircum{}bullet)|\_Z. Supplies the identification of the complex representing R c\_*(K|\_E) used to write the map of complexes.
\item Exactness of the pushforward along a finite morphism. Same chapter, the vanishing of the higher direct images of a finite morphism, used for the closed immersion i. Supplies that i\_*(F\textasciicircum{}bullet|\_Z) represents R i\_*(K|\_Z).
\end{itemize}\medskip
\ifdefined\hypertarget\hypertarget{sourcepointer-1}{}\fi
\noindent\textit{Source reference.} Original section 014Z: sources/stacks-chow/tags.txt, source lines 1490--1490. Prerequisite pointer invoked at proof line(s) 23368; it names a standard fact of a sibling chapter that the counted argument uses, and every step of the counted proof built on it is retained verbatim in the delivered body.\par\medskip
\ifdefined\hypertarget\hypertarget{sourcepointer-2}{}\fi
\noindent\textit{Source reference.} Original lemma 09V5: sources/stacks-etale-cohomology/pinned-original-source.tex, source lines 18708--18708. Prerequisite pointer invoked at proof line(s) 23348; it names a standard fact of the same chapter that the counted argument uses, and every step of the counted proof built on it is retained verbatim in the delivered body.\par\medskip
\ifdefined\hypertarget\hypertarget{sourcepointer-3}{}\fi
\noindent\textit{Source reference.} Original lemma 0A3U: sources/stacks-etale-cohomology/pinned-original-source.tex, source lines 18282--18282. Prerequisite pointer invoked at proof line(s) 23347; it names a standard fact of the same chapter that the counted argument uses, and every step of the counted proof built on it is retained verbatim in the delivered body.\par\medskip
\ifdefined\hypertarget\hypertarget{sourcepointer-4}{}\fi
\noindent\textit{Source reference.} Original proposition 03QP: sources/stacks-etale-cohomology/pinned-original-source.tex, source lines 8245--8245. Prerequisite pointer invoked at proof line(s) 23339; it names a standard fact of the same chapter that the counted argument uses, and every step of the counted proof built on it is retained verbatim in the delivered body.\par\medskip
\begin{lemma}
\label{lemma-blow-up-square-cohomological-descent}
Let $X$ be a scheme and let $Z \subset X$ be a closed subscheme
cut out by a quasi-coherent ideal of finite type. Consider the
corresponding blow up square
$$
\xymatrix{
E \ar[d]_\pi \ar[r]_j & X' \ar[d]^b \\
Z \ar[r]^i & X
}
$$
For $K \in D^+(X_\etale)$ with torsion cohomology sheaves
we have a distinguished triangle
$$
K \to Ri_*(K|_Z) \oplus Rb_*(K|_{X'}) \to Rc_*(K|_E) \to K[1]
$$
in $D(X_\etale)$ where $c = i \circ \pi = b \circ j$.
\end{lemma}

\begin{proof}
The notation $K|_{X'}$ stands for $b_{small}^{-1}K$.
Choose a bounded below complex $\mathcal{F}^\bullet$
of abelian sheaves representing $K$. Observe that
$i_*(\mathcal{F}^\bullet|_Z)$ represents $Ri_*(K|_Z)$
because $i_*$ is exact
(Proposition \ref{proposition-finite-higher-direct-image-zero}).
Choose a quasi-isomorphism
$b_{small}^{-1}\mathcal{F}^\bullet \to \mathcal{I}^\bullet$
where $\mathcal{I}^\bullet$ is a bounded below complex of injective
abelian sheaves on $X'_\etale$. This map is adjoint to a map
$\mathcal{F}^\bullet \to b_*(\mathcal{I}^\bullet)$ and
$b_*(\mathcal{I}^\bullet)$ represents $Rb_*(K|_{X'})$.
We have $\pi_*(\mathcal{I}^\bullet|_E) = (b_*\mathcal{I}^\bullet)|_Z$
by Lemma \ref{lemma-proper-base-change-f-star} and by
Lemma \ref{lemma-proper-base-change} this complex represents
$R\pi_*(K|_E)$. Hence the map
$$
Ri_*(K|_Z) \oplus Rb_*(K|_{X'}) \to Rc_*(K|_E)
$$
is represented by the surjective map of bounded below complexes
$$
i_*(\mathcal{F}^\bullet|_Z) \oplus
b_*(\mathcal{I}^\bullet)
\to
i_*\left(b_*(\mathcal{I}^\bullet)|_Z\right)
$$
To get our distinguished triangle it suffices to show that
the canonical map
$\mathcal{F}^\bullet \to i_*(\mathcal{F}^\bullet|_Z) \oplus
b_*(\mathcal{I}^\bullet)$
maps quasi-isomorphically onto the kernel of the map
of complexes displayed above (namely a short exact sequence
of complexes determines a distinguished triangle in the derived
category, see
Derived Categories, Section \ref{derived-section-canonical-delta-functor}).
We may check this on stalks at a geometric point $\overline{x}$ of $X$.
If $\overline{x}$ is not in $Z$, then $X' \to X$ is an isomorphism
over an open neighbourhood of $\overline{x}$. Thus, if $\overline{x}'$
denotes the corresponding geometric point of $X'$ in this case, then
we have to show that
$$
\mathcal{F}^\bullet_{\overline{x}} \to \mathcal{I}^\bullet_{\overline{x}'}
$$
is a quasi-isomorphism. This is true by our choice of $\mathcal{I}^\bullet$.
If $\overline{x}$ is in $Z$, then
$b_(\mathcal{I}^\bullet)_{\overline{x}} \to 
i_*\left(b_*(\mathcal{I}^\bullet)|_Z\right)_{\overline{x}}$
is an isomorphism of complexes of abelian groups. Hence the
kernel is equal to
$i_*(\mathcal{F}^\bullet|_Z)_{\overline{x}} =
\mathcal{F}^\bullet_{\overline{x}}$ as desired.
\end{proof}

\end{document}
